% generated by tools/render_paper_tables.py -- do not edit
\begin{description}[leftmargin=1.2em,style=nextline,font=\ttfamily\small]
\item[N\_modules] The pipeline is spread over N source files, none of them dominant; contrast \code{single\_script}.
\item[N\_runnable\_scripts] N files are runnable as programs. More than one means the documented single command is really a sequence.
\item[adaptive\_in\_name\_only] Every stated criterion names nothing but randomness: the round was labelled adaptive and was not.
\item[agree] The criterion stated in the log/report is the one the code computes.
\item[chain\_agreement] Share of a cell's replicates that chose the modal chain of methods — method reproducibility, which is not the same as score reproducibility.
\item[criterion\_switched] The acquisition criterion CHANGED between rounds — logic reacting to what it measured, rather than one fixed rule executed n times.
\item[deep\_ensemble] Several independently initialised models, averaged.
\item[enforced] A pilot runs and its success rate is compared against a threshold before the campaign proceeds — the task's 50\% pilot gate.
\item[evidence\_grounded] The round's validation error against baseline was recorded, so its choice can be checked against what was known at the time. An ungrounded round's reasoning is unfalsifiable.
\item[feasibility] Candidate choice weighted by whether solves there are expected to converge, steering away from regions with failed solves.
\item[gradient\_sensitivity] Points where the response is steep or curving — a sensitivity or Jacobian-driven criterion.
\item[hdf5] One HDF5 file holding many samples. Scales best and carries attributes/provenance, at the price of concurrent-write care.
\item[imported\_not\_subprocessed] predict.py is referenced but only imported or described, never run as the documented command. Code that works in the session's process and fails in a clean one is exactly what this misses.
\item[in\_process\_serial] Solves run one after another inside the driving process. Simplest, and the slowest: no parallelism, and one bad solve can poison the shared solver state.
\item[index] An index/manifest (CSV or JSONL) sits beside the arrays, listing every sample and its status. What makes a campaign resumable and auditable rather than a directory to be re-scanned.
\item[kernel\_ridge] Kernel ridge regression — a GP's mean without its variance.
\item[lcfs\_polygon] A grid point is inside the plasma iff it lies within the LCFS polygon — point-in-polygon against the D-shape boundary, whatever the helper is called (contains\_points, \_points\_in\_poly, inside\_lcfs\_mask). Geometric, per sample, and independent of the surrogate's own output.
\item[lhs] Latin hypercube: stratified in every dimension at once.
\item[mc\_dropout] Dropout left active at inference and sampled repeatedly — an ensemble's spread at one model's training cost.
\item[mismatch] The stated criterion and the implemented one have nothing in common. The strongest available signal that a run's account of itself is wrong.
\item[mlp\_torch] A hand-written PyTorch MLP. The most common choice in this corpus, usually mapping 5 parameters to PCA coefficients.
\item[model\_informed] The function that chooses points actually calls the surrogate. A model-derived criterion that never does is not one, whatever it is named.
\item[npz\_per\_solve] One compressed .npz per solved sample.
\item[oft\_gs\_domain] Meshed through OFT's own API (gs\_Domain / define\_region / build\_mesh) as the worked examples do. What the task's meshing milestone asks for, and what TokaMaker's topology checks accept.
\item[partial] Stated and implemented criteria overlap but do not match: one side names a component the other never does — see claimed-but-not-computed.
\item[pickle] Python-pickled objects (pickle/joblib/torch.save) — usually the model artifact rather than the dataset.
\item[prose\_only] A method named in the README or report.json that the code never evidences.
\item[random] Points drawn uniformly at random. As the CANDIDATE POOL this is normal and harmless; as the selection criterion it means the round was adaptive in name only.
\item[reload\_and\_check\_finite] A function both re-loads a written artifact and checks its finite fraction — the task's storage-validation gate, actually implemented. This is the check that catches an all-NaN grid stored as a success.
\item[residual\_ucb] Points are targeted where the model is MEASURABLY wrong — an error model fit on residuals or out-of-fold error, often with a UCB trade-off between exploiting known weakness and exploring. Uses measurement rather than the model's self-assessment, which can be confidently wrong.
\item[round\_cap] Stop because the 3-round cap was reached.
\item[single\_script] The whole pipeline is one or two files.
\item[singleton\_reused] One OFT\_env is built once and handed out on every call — a cached module global, a class guarding construction, or an lru\_cache. Honours OFT's one-env-per-process limit in the simplest way, at the cost of keeping every solve in one process: a solver crash takes the campaign with it.
\item[sobol] A Sobol low-discrepancy sequence — quasi-random, extensible.
\item[solver\_field\_eval] psi is read on the frozen 64x96 grid through the solver's own finite-element interpolator (get\_field\_eval). Evaluates the FE basis directly, so no second interpolation error is introduced.
\item[space\_filling] Purely geometric coverage of the input box: farthest-point/maximin, distance to the existing training set, k-means. Needs no model at all — which makes it robust, and makes it NOT adaptive in the sense of reacting to what was learned.
\item[stop\_threshold\_met] The campaign ended with the task's 70\% criterion satisfied — the intended way to finish early.
\item[stop\_unexplained] The campaign ended and no per-round validation error was recorded, so why it stopped cannot be read from the artifacts at all.
\item[stop\_without\_threshold] The campaign ended although validation error had NOT reached 0.30x baseline: the round cap, a budget, or a timeout ended it, not the stopping rule.
\item[subprocess\_per\_batch] The driver shells out (usually to its own script) to run a batch. Insulates solver state between batches and survives a hard crash; pays process start-up and serialises through files.
\item[subprocess\_reruns\_predict] Something in the workspace invokes predict.py as a subprocess with the contract CLI — the task's deliverable self-test, in a fresh process, which is the only way to catch an import-order or global-state dependency.
\item[test\_absent\_from\_fitting\_functions] No function that calls .fit()/backward()/optimizer.step() references a test path or array. Necessary, not sufficient: the task also forbids the test set influencing model SELECTION and ACQUISITION, which this check cannot see.
\item[top\_k] The highest-scoring candidates are taken (argsort/topk).
\item[training\_valid\_mask] The predictor masks with a valid-pixel mask derived from the training data (e.g. pixels finite in training). Cheap and stable, but the mask cannot adapt to a geometry unlike those seen in training — the same family of defect as the run that passed 9/9 gates while predicting plasma in vacuum.
\item[uncertainty] Points are ranked by the model's own predictive spread, without the record saying where that spread comes from. The family term; the two entries below are its specific forms.
\item[uncertainty\_ensemble] Spread across an ENSEMBLE of models (deep ensemble, bootstrap, or MC-dropout samples) — 'query by committee'. Needs no probabilistic model, and its quality depends entirely on the members disagreeing for real rather than sharing an initialisation.
\item[uncertainty\_gp] Posterior variance of a Gaussian process (or an acquisition built on it, e.g. expected improvement). Principled and calibrated where the GP's kernel assumptions hold; costly as the dataset grows.
\item[uniform\_random] Independent uniform draws, with no stratification.
\item[unverifiable\_from\_code] No function that chooses points could be classified, so the stated criterion can be neither confirmed nor contradicted.
\item[val\_threshold\_70pct] The task's own rule: stop once validation error is at most 0.30x the mean-predictor baseline, i.e. a 70\% reduction.
\end{description}
